\section{Data and method}
\label{sec:method}

\subsection{Data}
\label{sec:data}

We used four public keystroke datasets from three sources (Table~\ref{tab:data}). Two are editions of an
introductory programming (CS1) course at a US university, in which students wrote Python in a browser-based editor and
every edit was logged in the same format \cite{edwards2023cs1}; the 2019 and 2021 editions share the course, the
editor and the recording. KLiCKe is a corpus of argumentative essays written in one session each, with recorded press
times and a holistic score \cite{tian2025klicke,inputlog2026community}. CoAuthor is a dataset of creative and
argumentative writing with suggestions from a language model \cite{lee2022coauthor}; it is writing outside an
educational setting and enters as a contrast with a different recording, not as an educational sample. The method
needs only the time of each key press and whether it inserted or deleted a character, so it applies to any log with
these two fields. CS1 records edit events rather than key presses; an insertion of one character became a character
press and a deletion a deletion press by the mapping of the companion study \cite{miura2026runup} (S2).

A writer is a student (CS1, KLiCKe) or a participant (CoAuthor). We excluded writers with less than 10 minutes of
observation (47 CS1 2019 assignments, 84 KLiCKe essays and 2 CoAuthor sessions) and, in the count analyses, units with
less than 5 minutes of exposure. Two random holdouts set aside in earlier work had been opened in the companion study
and are included (CS1 2019 and KLiCKe); the holdouts of CS1 2021 and CoAuthor were not opened. KLiCKe is the only
dataset with an outcome that we used: in the 2019 CS1 release the grade table was withdrawn by the data authors as
erroneous, and the 2021 grades were not used because only 31 writers were on the seen side
(Section~\ref{sec:limitations}).

\begin{table}[t]
\centering\footnotesize
\caption{The four datasets. Units are assignments (CS1), essays (KLiCKe) and sessions (CoAuthor). Onsets are the
stall onsets under the baseline definition. Holdout: writers set aside before any analysis.}
\label{tab:data}
\begin{tabularx}{\textwidth}{@{}l X r r r X X@{}}
\toprule
Dataset & Writers and setting; one event & Writers & Units & Onsets & Outcome used & Holdout \\
\midrule
CS1 2019 & students of a CS1 course writing Python; one logged edit & 459 & 1,469 & 33,710 & none (grade table
withdrawn) & 146 writers opened in the companion study, included \\
CS1 2021 & same course, editor and recording & 31 & 177 & 6,217 & none (31 writers) & 13 writers, not opened \\
KLiCKe & students writing one argumentative essay; one key press & 4,908 & 4,908 & 33,003 & holistic score (1--6)
& 1,498 writers opened earlier, included \\
CoAuthor & adults writing with model suggestions; one key press & 41 & 1,155 & 1,518 & none & 18 writers
(288 sessions), not opened \\
\bottomrule
\end{tabularx}
\end{table}

\subsection{The baseline definition and the detector}
\label{sec:detector}

A stall is detected from three ingredients (Figure~\ref{fig:detector}). Let \emph{condition A} hold at time $t$ when,
in the window $(t - W, t]$, the number of character presses minus the number of deletion presses is less than one. When
condition A rises from false to true at time $s$ and stays true until $t_0 = s + \theta$, an onset is raised at $t_0$.
The stall lasts from $s$ to the next character press in the same observation interval (the resumption); its length is
measured from $t_0$. The \emph{baseline definition} is the default of the detector's reference implementation: $W =
30$~s, $\theta$ the 0.99 nearest-rank quantile of the writer's inter-keystroke intervals so far in the unit (defined
once 100 intervals exist), the net rule for condition A, and the default observation intervals (the spans in which the
estimator, re-run on the recorded stream, wrote its per-second rows, split where two rows are more than 2~s apart).
It is the definition of the companion study's appendix and was fixed before any KLiCKe score was read. The label marks a pause in character
production only (Section~\ref{sec:limitations}). The detector reproduced the
onsets of the reference implementation in every unit of the four datasets (S2).

\begin{figure}[t]
\centering
\begin{tikzpicture}[x=0.11cm,y=0.5cm,font=\footnotesize]
  \draw[->] (0,0) -- (120,0) node[right] {time};
  \foreach \x in {2,5,7,10,13,15,18,21,24,27} \draw (\x,0) -- (\x,0.45);
  \foreach \x in {96,99,103,106,110} \draw (\x,0) -- (\x,0.45);
  \node[above] at (14,0.5) {character presses};
  \draw[fill=figband,draw=none] (27,-0.35) rectangle (57,-0.95);
  \node at (42,-0.65) {$W$};
  \draw[thick] (57,0.3) -- (57,-2.2) node[below] {$s$};
  \draw[thick] (81,0.3) -- (81,-2.2) node[below] {$t_0$};
  \draw[<->] (57,-1.55) -- node[above] {$\theta$} (81,-1.55);
  \draw[thick] (96,0.3) -- (96,-2.2) node[below] {resumption};
  \draw[decorate,decoration={brace,mirror,amplitude=4pt}] (57,-3.3) -- node[below=4pt] {stall (length measured
  from $t_0$)} (96,-3.3);
\end{tikzpicture}
\caption{The detector. Condition A (fewer than one net character press in the window $W$, shaded) becomes true at $s$; if it
stays true for $\theta$, an onset is raised at $t_0 = s + \theta$. The stall ends at the next character press.}
\label{fig:detector}
\end{figure}

\subsection{The four axes and the 89 definitions}
\label{sec:grid}

We varied four choices that every such detector has to make, each over levels drawn from prior practice or set around
the default (Table~\ref{tab:grid}). The full crossing gives 90 definitions; with the baseline removed, 89 remain, and
all proportions below are out of 89. The two families of $\theta$ are different constructs: a quantile of the writer's
own intervals marks a pause that is long \emph{for this writer}, a fixed number of seconds marks a pause that is long in
absolute terms. We report them together because both are in use, and separately where they part
(53 definitions with a personal quantile, 36 with fixed seconds).

\begin{table}[t]
\centering\footnotesize
\caption{The four axes. The baseline level is marked with an asterisk.}
\label{tab:grid}
\begin{tabularx}{\textwidth}{@{}l l X@{}}
\toprule
Axis & Levels & Source of the levels \\
\midrule
Window $W$ & 15, 30$^\ast$, 60~s & 30~s is the default; 15 and 60~s are half and double the default (we found no prior
use of a 15-s window) \\
Threshold $\theta$ & writer quantile 0.95, 0.99$^\ast$, 0.995; fixed 10, 30~s & writer-relative thresholds
\cite{chukharev2014pauses}; fixed thresholds in writing span about 0.25--10~s \cite{torrance2024methods}, and 30~s equals
the default window \\
Condition A & net$^\ast$ (characters $-$ deletions $< 1$); characters only ($< 1$) & whether deletions cancel character
presses \\
Observation & default intervals$^\ast$; split at key-press gaps of 5~min; of 30~min & 5~min of inactivity as
time-on-task \cite{hart2023time}; 30~min as a session timeout \cite{halfaker2015session} \\
\bottomrule
\end{tabularx}
\end{table}

\subsection{What we measured}
\label{sec:measures}

\paragraph{Rates and ratios.} The stall rate is the number of onsets per observed hour, pooled over a dataset's units.
For each definition we compared the rank order of the four datasets with the baseline order, and the ratio of the rates
of each pair of datasets with the baseline ratio.

\paragraph{The essay-score association.} For KLiCKe, the Spearman correlation between a writer's stall rate and the
holistic score, and the median of the correlations within each of the four prompts. Because the sample is large, weak
correlations have intervals that exclude zero, so we also read magnitudes, using $|\rho| < 0.1$ as negligible and 0.3
as moderate \cite{cohen1988power} [D].

\paragraph{Typing speed} [E]. The median of a writer's inter-keystroke intervals shorter than 2~s, a proxy for typing
speed that does not contain the stall labels. We computed the partial Spearman correlation of score and stall rate
given speed, with six ways of measuring speed (S3).

\paragraph{Overdispersion within writers.} For sub-writer units (CS1 work blocks and assignments; CoAuthor sessions),
the Pearson dispersion index $\varphi = \sum (n - \hat e)^2 / \hat e \,/\, (k - p)$, where the expected count $\hat e$ is
the writer's own rate times the unit's exposure and $p$ the number of writers. $\varphi = 1$ is what counting at a
constant rate within a writer would give; 1.5 is the value at which we would switch a count model from Poisson to
negative binomial (S4).

\paragraph{Components of stall length.} A linear mixed model of log stall length with random effects of writer,
assignment, day and work block (CS1; days are runs of blocks separated by gaps of 6~h or more) or of writer and
session (CoAuthor), fitted by REML; we report the share of each level and of the residual, the stall-to-stall
variation. The fit took the maximum of the REML likelihood over three optimisers (S5). We do not report the analogous
decomposition of counts, whose variational estimates did not recover simulated components and moved between runs on
the same input (S6).

\paragraph{Intervals.} All intervals are 95\% percentile intervals from 2,000 resamples of writers. The writer quantile
$\theta$ is computed within each unit and reused in the resample, so its estimation error is not in the interval
(S2). The interval of a definition shows sampling variation; the spread across definitions shows definitional
variation, and we set the two side by side \cite{robitzsch2022multiverse}. Definitions share axes and are not
independent, so we read counts as proportions of the grid and do not test the grid as a whole
\cite{simonsohn2020specification}.

\subsection{Status of the analyses}
\label{sec:status}

Every computation ran under conditions committed in the author's repository before it ran; these are not third-party
registrations (S1). We mark each result with its status: [P] primary, with the quantities and criteria committed
before any result; [A] added after the primary results, with the quantity and its criterion committed before that
quantity was seen (the length decomposition over all definitions, the intervals and the score association); [E]
exploratory, chosen after seeing results; [D] descriptive, counted after seeing results. Three changes bear on the
claims. First, the covariate for the score association was chosen twice after seeing results, from the key-press rate
to typing speed and then to six ways of measuring speed (S1); this is why the adjusted results are [E]. Second, the pair
we had committed to report first, the two CS1 editions, did not reach the criterion for a reversal of the rate ratio
(Section~\ref{sec:res-ratio}). Third, the criterion for a reversal of the score association required only one definition
with an interval on the opposite side of zero, which is lenient; we therefore report the full distribution. The
baseline definition, the conversion, the estimator and the opened holdouts are shared with the companion study, which
tested a different question on part of the same data (S9).
